\section{System overview}

% claims: C01 C04 C05 C13; sources: S01 S02 S05 S06 S13
Our construction joins four layers that are useful separately but answer different
questions together.  A \emph{causal computation} produces an internal state that
the answer must use.  \emph{Receipts} authenticate what trusted instrumentation
observed.  A verified state class selects a \emph{carrier authority}: a keyed
identity that changes token probabilities.  Finally, an output-side detector asks
which authority the generated text statistically supports.  We call the resulting
regularity a \emph{carrier}.  It is evidence of compatibility with a state class,
not an independent proof of the execution that produced the text.

\begin{figure}[t]
  \centering
  \small
  \fbox{\shortstack{1. Causal computation\\state affects answer}}
  $\rightarrow$
  \fbox{\shortstack{2. Receipts\\authenticate events}}
  $\rightarrow$
  \fbox{\shortstack{3. State authority\\biases sampling}}
  $\rightarrow$
  \fbox{\shortstack{4. Output verification\\scores authorities}}
  % claims: C01 C02 C04 C05; sources: S03 S04 S05 S06 S13
  \caption{The four-layer system. Computation establishes causal use; exact,
  abstract, and chain receipts establish trusted provenance claims; a
  state-derived authority controls a statistical signal; and output scoring
  recovers that signal. No single layer establishes all four properties.}
  \label{fig:overview}
\end{figure}

\paragraph{What each receipt means.}
% claims: C04 C05; sources: S05 S06 S13
An \emph{exact receipt} identifies one particular internal event in one
execution.  An \emph{abstract receipt} says that the event belongs to a
registered state class, so two different executions may share the same abstract
state while having different exact receipts.  A \emph{causal-chain receipt}
authenticates the registered order and relation from \ztwo{} through \zthree{} to
the answer.  These distinctions matter: a different execution can fail exact
provenance while remaining compatible with the same abstract state and hence
the same carrier authority.

\paragraph{A running matched example.}
% claims: C03 C09; sources: S04 S10 S12 S13
For input $(4,15,1,1)$, the natural execution produces
$(\ztwo,\zthree,y)=(4,5,5)$.  We then intervene on the internal state to obtain
an alias execution with $(12,13,5)$.  The internal path changes but the answer
does not.  The abstract receipt for state 4 selects authority $A$; the receipt
for state 12 selects authority $B$.  Both executions reuse the same prompt,
answer, semantic payload, configuration, output indices, and indexed random
draws.  The carrier is therefore expected to switch from $A$ to $B$ only
because the authenticated causal state switched.

\begin{figure}[t]
  \centering
  \small
  \begin{tabularx}{0.98\linewidth}{@{}YcY@{}}
    \toprule
    \textbf{Natural execution} & \textbf{Held fixed} & \textbf{Alias execution} \\
    \midrule
    $\ztwo=4\rightarrow\zthree=5\rightarrow y=5$
      & prompt; answer; payload; configuration; indexed draws
      & $\ztwo=12\rightarrow\zthree=13\rightarrow y=5$ \\
    abstract state 4 $\rightarrow$ authority $A$
      & same visible task and semantic content
      & abstract state 12 $\rightarrow$ authority $B$ \\
    output carries signal $A$
      & only intended switching source: authenticated \ztwo{}
      & output carries signal $B$ \\
    \bottomrule
  \end{tabularx}
  % claims: C03 C09; sources: S04 S10 S12 S13
  \caption{The decisive natural/alias contrast. The visible problem, final
  answer, semantic payload, generation settings, and position-indexed random
  draws are paired. The authenticated \ztwo{} state, downstream \zthree{},
  receipt context, derived authority, and intended carrier signal change.}
  \label{fig:contrast}
\end{figure}

\paragraph{Experimental unit.}
% claims: C08; sources: S10 S11
A \emph{closure} is one complete matched natural/alias test unit, including its
two executions, receipts, generated outputs, controls, and frozen verdicts.  It
passes only if both executions satisfy every registered computation, receipt,
chain, parsing, carrier, and attack gate.  Outputs, token positions, authority
scores, and attack checks within a closure are dependent measurements, not
independent replications.
